\chapter{Bài toán và ý nghĩa thực tế}\label{chap:problem}
\section{Động cơ thực hiện}
Trong tích hợp dữ liệu, cùng một nội dung bệnh học có thể được mô tả bằng
các hệ thống mã khác nhau. Một bảng ánh xạ cung cấp các cặp mã ứng viên,
nhưng việc có mặt trong bảng chưa giải quyết được câu hỏi về mức độ phù hợp
của từng cặp. Người kiểm tra cần biết nhãn có tương ứng hay không, có khác
mức độ chi tiết hay có từ ngữ chỉ tình trạng xung đột. Nhu cầu hỗ trợ chuyên
gia kiểm tra các ứng viên ánh xạ cũng là bối cảnh của công trình
OntologyRAG~\cite{ontologyrag2025}.

Đồ án sử dụng mã nguồn và dữ liệu hiện có trong repository làm nền tảng,
tập trung vào lớp tri thức COKB. Vấn đề cốt lõi là xây dựng một quy trình
\emph{đánh giá có căn cứ}: cùng dữ kiện và cùng phiên bản luật phải cho
cùng nhận định, đồng thời cho phép truy ngược nhận định đó đến tiền đề và nguồn.

\section{Phát biểu bài toán}
Mỗi ứng viên ánh xạ được mô hình hóa bởi
\begin{equation}
  m=(id,c_s,c_t,e,b),\qquad c=(s,k,\ell),
  \label{eq:mapping}
\end{equation}
trong đó $id$ là định danh; $c_s,c_t$ là khái niệm nguồn và đích;
$s$ là tên hệ mã; $k$ là giá trị mã; $\ell$ là nhãn;
$e$ là nguồn bằng chứng; $b$ là tác nhân khai báo.
Đầu ra của hệ thống là
\begin{equation}
  f_{\K}(m)=(y,\tau,\pi),\qquad
  y\in\{A,B,C,\bot\},\quad \tau=\texttt{rule\_based},
  \label{eq:task}
\end{equation}
với $\bot$ tương ứng \Unassessed{} và $\pi$ là vết suy luận.
Tập tri thức $\K$ bao gồm mô hình dữ liệu, các phép xử lý và tập luật được
chọn cho lần chạy. Việc tìm ra mã đích mới ngoài các ứng viên đầu vào
không nằm trong hàm $f_{\K}$.

\begin{table}[htbp]
\centering\small
\caption{Ý nghĩa đầu ra và cách sử dụng dự kiến}\label{tab:practical}
\begin{tabularx}{\textwidth}{lYY}
\toprule
Đầu ra & Diễn giải theo bộ luật hiện tại & Hỗ trợ công việc nào?\\
\midrule
A & Nhãn trùng nhau sau chuẩn hóa. & Đưa các cặp dễ kiểm tra vào nhóm đề xuất xác nhận.\\
B & Thỏa điều kiện bao hàm từ hoặc khái quát hóa theo heuristic. & Chú ý khác biệt về nhãn; kiểm tra thêm ngữ cảnh.\\
C & Phát hiện cặp từ/cụm từ đối lập đã khai báo. & Đưa vào nhóm cảnh báo xung đột.\\
$\bot$ & Không có luật A/B/C nào khớp. & Giữ lại cho chuyên gia hoặc bổ sung bằng chứng.\\
\bottomrule
\end{tabularx}
\end{table}

Các hành động nghiệp vụ trong Bảng~\ref{tab:practical} là cách sử dụng
được đề xuất. Code hiện trả assessment và proof; chưa có bộ chính sách
phê duyệt, hàng đợi chuyên gia hay chức năng tự động chuyển dữ liệu bệnh viện.
Đặc biệt, A là kết quả so sánh nhãn theo thuật toán, không phải chứng nhận
tương đương y khoa giữa hai khái niệm.

\section{COKB mang lại giá trị gì?}
Giá trị thứ nhất là \textbf{biến tiêu chí kiểm tra thành tri thức có thể thực thi}.
Ví dụ, tiêu chí “hai nhãn trùng nhau sau chuẩn hóa” được đóng gói thành
một luật có định danh. Khi tiêu chí thay đổi, người phát triển biết phải
sửa thành phần nào và đo lại ảnh hưởng lên dữ liệu nào.

Giá trị thứ hai là \textbf{giải trình được một quyết định cụ thể}.
Với một cảnh báo xung đột, hệ thống cung cấp cả hai nhãn và danh sách
cụm từ đã kích hoạt luật. Chuyên gia có thể nhận ra cảnh báo đúng hoặc phát
hiện lỗi do thuật toán hiểu sai phủ định. Đây là lợi ích trực tiếp của
proof trace đối với kiểm tra chất lượng dữ liệu.

Giá trị thứ ba là \textbf{biểu diễn rõ giới hạn hiểu biết}.
\Unassessed{} giữ lại những cặp mà luật chưa xử lý được. Nếu được tích hợp
vào quy trình rà soát, trạng thái này giúp hình thành danh sách cần kiểm tra
và giúp nhóm phát triển ưu tiên bổ sung tri thức. Tác dụng tiết kiệm thời gian
cho chuyên gia vẫn cần được đo bằng nghiên cứu người dùng; đồ án chưa có
thực nghiệm xác nhận lợi ích định lượng đó.

\section{Mục tiêu và phạm vi}
Đồ án hướng đến bốn kết quả: mô hình hóa các đối tượng và quan hệ của ánh xạ;
hiện thực các phép kiểm tra cùng luật đánh giá; xuất được kết luận kèm vết
suy luận; và đánh giá chất lượng bằng dữ liệu có nhãn tham chiếu.
Các ví dụ ICD chỉ được diễn giải theo nhãn có trong repository.

Phạm vi thực nghiệm gồm năm mapping mẫu và 500 cặp nhãn gold.
Hệ thống chưa xử lý đầy đủ ngữ nghĩa phân cấp bệnh, phiên bản ICD,
ngữ cảnh lâm sàng hoặc ánh xạ đích dạng tổ hợp. Báo cáo phân biệt các
khả năng đã chạy được, các khái niệm chỉ mới khai báo trong ontology
và các hướng ứng dụng cần triển khai thêm.
